\documentclass[12pt]{article}

% Packages
\usepackage[margin=2.5cm]{geometry}
\usepackage{amsmath,amssymb,amsthm}
\usepackage{bm}
\usepackage{enumitem}
\usepackage{booktabs}
\usepackage{graphicx}
\usepackage[hidelinks]{hyperref}
\usepackage{listings}
\usepackage{xcolor}

% Code listing style
\lstset{
  language=Matlab,
  basicstyle=\small\ttfamily,
  keywordstyle=\color{blue}\bfseries,
  commentstyle=\color{green!50!black}\itshape,
  stringstyle=\color{red!70!black},
  numbers=left,
  numberstyle=\tiny\color{gray},
  numbersep=5pt,
  frame=single,
  breaklines=true,
  breakatwhitespace=true,
  tabsize=4,
  showstringspaces=false,
  captionpos=b,
}

% Theorem environments
\newtheorem{proposition}{Proposition}
\newtheorem{theorem}{Theorem}
\newtheorem{corollary}{Corollary}
\newtheorem{remark}{Remark}

% No indentation
\setlength{\parindent}{0pt}
\setlength{\parskip}{6pt}

\begin{document}

\title{Universal Route for Adsorption Isotherm Analysis:\\
Convergence to Classical Models and Fitting Implementation\\[6pt]
\large (Version 5)}
\author{}
\date{}
\maketitle

% ======================================================================
\section{Introduction}
% ======================================================================

In the preceding versions (v2--v3) we established a model-agnostic
protocol for extracting $Q_{\max}$ and
$K_{\mathrm{aff}}\equiv K_H/Q_{\max}$ from adsorption data, based on
the heterogeneous Langmuir integral
\begin{equation}
  q(C) = K_p\,C + \int_{0}^{\infty} \rho(K)\,\frac{KC}{1+KC}\,dK,
  \qquad \rho(K)\ge 0,
  \label{eq:master}
\end{equation}
with invariants
\begin{equation}
  Q_{\max} = \int_0^\infty \rho(K)\,dK, \qquad
  K_H = \int_0^\infty K\,\rho(K)\,dK, \qquad
  K_{\mathrm{aff}} = \frac{K_H}{Q_{\max}}.
  \label{eq:invariants}
\end{equation}

The purpose of this document is to \emph{prove} that every classical
adsorption isotherm that satisfies assumptions A1--A4 is a special
case of Eq.~\eqref{eq:master}, and that the universal invariants
$Q_{\max}$, $K_H$, and $K_{\mathrm{aff}}$ reduce exactly to the
known parameters of each model.  We also characterise the degenerate
cases (Freundlich, Redlich--Peterson with $\beta<1$) where one or more
invariants diverge, and prove a kernel-independence theorem that makes
the invariants robust to the choice of local isotherm mechanism.

% ======================================================================
\section{Framework recap}
\label{sec:framework}
% ======================================================================

\textbf{Assumptions} (from v2--v3):
\begin{itemize}[nosep]
  \item[A1.] Independent-site superposition (heterogeneous Langmuir
        class, no cooperativity).
  \item[A2.] Finite saturating capacity:
        $Q_{\max}=\int_0^\infty\rho(K)\,dK<\infty$.
  \item[A3.] Finite first moment:
        $K_H=\int_0^\infty K\rho(K)\,dK<\infty$.
  \item[A4.] Sufficient experimental coverage of the Henry and plateau
        regimes.
\end{itemize}

The \emph{normalised} affinity distribution is
$f(K)\equiv\rho(K)/Q_{\max}$, with $\int_0^\infty f(K)\,dK=1$.
The effective affinity is the first moment of $f$:
\begin{equation}
  K_{\mathrm{aff}} = \int_0^\infty K\,f(K)\,dK = \langle K\rangle_f.
  \label{eq:Kaff_moment}
\end{equation}

% ======================================================================
\section{Exact convergence to classical models}
\label{sec:convergence}
% ======================================================================

% ------ 3.1 Langmuir ------
\subsection{Langmuir isotherm}

\begin{proposition}[Langmuir]
\label{prop:langmuir}
Let $\rho(K)=Q_m\,\delta(K-K_0)$ with $Q_m>0$ and $K_0>0$.  Then:
\begin{enumerate}[nosep]
  \item $q(C) = \displaystyle\frac{Q_m K_0 C}{1+K_0 C}$
        \quad (single-site Langmuir),
  \item $Q_{\max}=Q_m$, \quad $K_H=Q_m K_0$, \quad
        $K_{\mathrm{aff}}=K_0$.
\end{enumerate}
\end{proposition}

\begin{proof}
Substituting $\rho(K)=Q_m\,\delta(K-K_0)$ into
Eq.~\eqref{eq:master} (with $K_p=0$):
\[
  q(C)=\int_0^\infty Q_m\,\delta(K-K_0)\,\frac{KC}{1+KC}\,dK
      = Q_m\,\frac{K_0 C}{1+K_0 C}.
\]
The invariants follow directly:
$Q_{\max}=\int Q_m\,\delta(K-K_0)\,dK=Q_m$,\;
$K_H=\int K\,Q_m\,\delta(K-K_0)\,dK=Q_m K_0$,\;
$K_{\mathrm{aff}}=K_0$.
All assumptions A1--A4 are satisfied.
\end{proof}

% ------ 3.2 Multi-site Langmuir ------
\subsection{Multi-site Langmuir isotherm}

\begin{proposition}[Multi-site Langmuir]
\label{prop:multilang}
Let $\rho(K)=\sum_{j=1}^{J}Q_j\,\delta(K-K_j)$ with
$Q_j>0$, $K_j>0$ distinct.  Then:
\begin{enumerate}[nosep]
  \item $q(C)=\displaystyle\sum_{j=1}^{J}
        \frac{Q_j K_j C}{1+K_j C}$
        \quad ($J$-site Langmuir),
  \item $Q_{\max}=\displaystyle\sum_{j=1}^{J}Q_j$, \quad
        $K_H=\displaystyle\sum_{j=1}^{J}Q_j K_j$, \quad
        $K_{\mathrm{aff}}=\displaystyle
        \frac{\sum_j Q_j K_j}{\sum_j Q_j}$.
\end{enumerate}
\end{proposition}

\begin{proof}
Linearity of the integral and the sifting property of the
Dirac delta give (1) immediately.  The invariants are finite sums
of positive terms, hence finite.  $K_{\mathrm{aff}}$ is the
capacity-weighted arithmetic mean of the individual affinities.
\end{proof}

\begin{corollary}[Bounds on $K_{\mathrm{aff}}$]
\label{cor:bounds}
For any discrete mixture,
$\min_j K_j \le K_{\mathrm{aff}} \le \max_j K_j$.
\end{corollary}

\begin{proof}
This follows from $K_{\mathrm{aff}}$ being a convex combination
of the $K_j$ with positive weights $Q_j/Q_{\max}$.
\end{proof}

% ------ 3.3 Sips ------
\subsection{Sips (Langmuir--Freundlich) isotherm}

\begin{proposition}[Sips]
\label{prop:sips}
Let the affinity distribution be the Sips
distribution~\cite{Sips1948}:
\begin{equation}
  f(K) = \frac{\sin(m\pi)}{\pi}\,
  \frac{(K/K_s)^{m}}
       {K\bigl[(K/K_s)^{2m}+2\cos(m\pi)(K/K_s)^{m}+1\bigr]},
  \qquad 0<m\le 1,\; K_s>0.
  \label{eq:sips_dist}
\end{equation}
Then:
\begin{enumerate}[nosep]
  \item $q(C) = \displaystyle\frac{Q_{\max}(K_s C)^m}
        {1+(K_s C)^m}$
        \quad (Sips isotherm).
  \item $Q_{\max}$ is finite and correctly recovered:
        $\lim_{C\to\infty}q(C)=Q_{\max}$.
  \item For $m=1$: the Sips distribution degenerates to
        $\delta(K-K_s)$, so $K_H=Q_{\max}K_s$ and
        $K_{\mathrm{aff}}=K_s$ (Langmuir limit).
  \item For $0<m<1$: $K_H=+\infty$ (first moment diverges).
        Assumption A3 is violated.
\end{enumerate}
\end{proposition}

\begin{proof}
\textit{Part (1):} The isotherm follows from the Stieltjes-transform
inversion by Sips~\cite{Sips1948}; see also
Jaroniec and Madey~\cite{Jaroniec1988}, \S3.2.

\textit{Part (2):} Since $(K_s C)^m/(1+(K_s C)^m)\to 1$ as
$C\to\infty$, the limit gives $Q_{\max}$.

\textit{Part (3):} As $m\to 1$, $\sin(m\pi)\to 0$ and the
distribution narrows; by L'H\^opital's rule the Sips distribution
converges to $\delta(K-K_s)$ in the distributional
sense~\cite{Jaroniec1988}.

\textit{Part (4):} For $0<m<1$, differentiate the isotherm:
\[
  \frac{dq}{dC}\bigg|_{C\to 0}
  = Q_{\max}\,m\,K_s^m\,\lim_{C\to 0^+}C^{m-1}
  = +\infty,
\]
since $m-1<0$.  Equivalently, the tail of $f(K)$ decays as
$f(K)\sim K^{-m-1}$ for $K\to\infty$, so
$\int_0^\infty K f(K)\,dK$ diverges.
\end{proof}

\begin{remark}[Operational $K_{\mathrm{aff}}$ for Sips]
Although $K_H=\infty$ for $m<1$, the universal workflow produces a
\emph{finite} operational $\widehat{K}_{\mathrm{aff}}$ from data
spanning a finite concentration range $[C_{\min},C_{\max}]$.  This
estimate depends on the range and should be reported together with
the concentration window used.  The sensitivity diagnostic
(Section~5.3 of v3) flags this case: $K_{\mathrm{aff}}$ varies
substantially with the size of the Henry region $\mathcal{I}_0$.
\end{remark}

% ------ 3.4 Tóth ------
\subsection{T\'oth isotherm}

\begin{proposition}[T\'oth]
\label{prop:toth}
The T\'oth isotherm~\cite{Toth1971}
\begin{equation}
  q(C) = \frac{Q_{\max}\,K_T\,C}{\bigl(1+(K_T C)^t\bigr)^{1/t}},
  \qquad 0<t\le 1,\; K_T>0,
  \label{eq:toth}
\end{equation}
satisfies the heterogeneous Langmuir integral with a specific
$\rho(K)$ related to one-sided L\'evy stable
distributions~\cite{Jaroniec1975}.  Its invariants are:
\begin{enumerate}[nosep]
  \item $Q_{\max}=Q_{\max}$ (finite for all $t$).
  \item $K_H = Q_{\max}\,K_T$ (finite for all $0<t\le 1$).
  \item $K_{\mathrm{aff}} = K_T$.
  \item For $t=1$: reduces to the Langmuir isotherm.
\end{enumerate}
\end{proposition}

\begin{proof}
\textit{Part (1):}
\[
  \lim_{C\to\infty} \frac{Q_{\max}K_T C}{(1+(K_T C)^t)^{1/t}}
  = \lim_{C\to\infty} \frac{Q_{\max}K_T C}{(K_T C)^{t/t}}
    \cdot\frac{1}{(1+(K_TC)^{-t})^{1/t}}
  = Q_{\max}\cdot 1 = Q_{\max}.
\]

\textit{Part (2):}  Differentiate at $C=0$.  Write
$q(C)=Q_{\max}K_T C\cdot g(C)$ where
$g(C)=(1+(K_T C)^t)^{-1/t}$.  Then $g(0)=1$ and
\[
  K_H = \frac{dq}{dC}\bigg|_{C=0} = Q_{\max}K_T\cdot g(0) = Q_{\max}K_T.
\]

\textit{Part (3):} $K_{\mathrm{aff}}=K_H/Q_{\max}=K_T$.

\textit{Part (4):} For $t=1$:
$(1+K_T C)^{1/1}=1+K_T C$, giving the standard Langmuir.
\end{proof}

\begin{remark}[Heterogeneity parameter $t$]
The T\'oth parameter $t$ controls asymmetry of the energy
distribution.  For $t<1$ the distribution is broader with a longer
tail toward weak sites ($K\to 0$).  Unlike the Sips parameter $m$,
the T\'oth parameter does \emph{not} cause the first moment to
diverge; hence $K_{\mathrm{aff}}=K_T$ is well-defined for all
$t\in(0,1]$.  This makes the T\'oth model particularly compatible
with the universal route.
\end{remark}

% ------ 3.5 Temkin ------
\subsection{Temkin isotherm}

\begin{proposition}[Temkin]
\label{prop:temkin}
Let the affinity distribution be uniform in $\ln K$\,:
\begin{equation}
  \rho(K) = \frac{Q_{\max}}{K\ln(K_{\max}/K_{\min})},
  \qquad K_{\min}\le K \le K_{\max},
  \label{eq:temkin_dist}
\end{equation}
with $0<K_{\min}<K_{\max}<\infty$.  Then:
\begin{enumerate}[nosep]
  \item $q(C) = \displaystyle\frac{Q_{\max}}{\ln(K_{\max}/K_{\min})}
        \bigl[\ln(1+K_{\max}C)-\ln(1+K_{\min}C)\bigr]$.
  \item $Q_{\max}=Q_{\max}$ (finite).
  \item $K_H = \displaystyle\frac{Q_{\max}(K_{\max}-K_{\min})}
        {\ln(K_{\max}/K_{\min})}$ (finite).
  \item $K_{\mathrm{aff}} = \displaystyle
        \frac{K_{\max}-K_{\min}}{\ln(K_{\max}/K_{\min})}$
        \quad (logarithmic mean of $K_{\min}$ and $K_{\max}$).
  \item In the intermediate range $K_{\min}C\ll 1\ll K_{\max}C$:
        $q\approx b\,\ln(K_{\max}C)$ with
        $b\equiv Q_{\max}/\ln(K_{\max}/K_{\min})$
        (classical Temkin form).
\end{enumerate}
\end{proposition}

\begin{proof}
\textit{Part (1):}
\begin{align*}
  q(C) &= \int_{K_{\min}}^{K_{\max}}
    \frac{Q_{\max}}{K\ln(K_{\max}/K_{\min})}
    \,\frac{KC}{1+KC}\,dK
  = \frac{Q_{\max}}{\ln(K_{\max}/K_{\min})}
    \int_{K_{\min}}^{K_{\max}}
    \frac{C}{1+KC}\,dK \\
  &= \frac{Q_{\max}}{\ln(K_{\max}/K_{\min})}
    \bigl[\ln(1+KC)\bigr]_{K_{\min}}^{K_{\max}}
  = \frac{Q_{\max}}{\ln(K_{\max}/K_{\min})}
    \bigl[\ln(1+K_{\max}C)-\ln(1+K_{\min}C)\bigr].
\end{align*}

\textit{Part (2):} As $C\to\infty$:
$\ln(1+K_{\max}C)-\ln(1+K_{\min}C)\to
\ln(K_{\max}C)-\ln(K_{\min}C)=\ln(K_{\max}/K_{\min})$,
so $q\to Q_{\max}$.

\textit{Part (3):}
\begin{align*}
  K_H &= \int_{K_{\min}}^{K_{\max}}
    K\cdot\frac{Q_{\max}}{K\ln(K_{\max}/K_{\min})}\,dK
  = \frac{Q_{\max}}{\ln(K_{\max}/K_{\min})}
    \int_{K_{\min}}^{K_{\max}} dK
  = \frac{Q_{\max}(K_{\max}-K_{\min})}
        {\ln(K_{\max}/K_{\min})}.
\end{align*}

\textit{Part (4):} Division of (3) by (2).

\textit{Part (5):} When $K_{\min}C\ll 1$,
$\ln(1+K_{\min}C)\approx K_{\min}C\approx 0$.
When $K_{\max}C\gg 1$,
$\ln(1+K_{\max}C)\approx \ln(K_{\max}C)$.
Hence $q\approx b\ln(K_{\max}C)$, which is the classical
Temkin equation $q=b\ln(AC)$ with $A=K_{\max}$.
\end{proof}

\begin{remark}[Logarithmic mean affinity]
The result $K_{\mathrm{aff}}=(K_{\max}-K_{\min})/
\ln(K_{\max}/K_{\min})$ is the \emph{logarithmic mean}
of $K_{\min}$ and $K_{\max}$, satisfying
$K_{\min}<K_{\mathrm{aff}}<K_{\max}$.
As $K_{\min}\to K_{\max}$ (homogeneous surface),
$K_{\mathrm{aff}}\to K_{\max}$ and the isotherm reduces to a
single-site Langmuir (Proposition~\ref{prop:langmuir}).
\end{remark}

% ------ 3.6 Dual-mode ------
\subsection{Dual-mode (partition + Langmuir) isotherm}

\begin{proposition}[Dual-mode]
\label{prop:dualmode}
Let $q(C)=K_p C + Q_m K_0 C/(1+K_0 C)$ with $K_p>0$, $Q_m>0$,
$K_0>0$.  Then:
\begin{enumerate}[nosep]
  \item The saturating component has
        $\rho(K)=Q_m\,\delta(K-K_0)$,
        and the linear component is $K_p C$.
  \item $Q_{\max}=Q_m$ (saturating capacity only).
  \item $K_H=Q_m K_0$ (Henry constant of the saturating part).
  \item $K_{\mathrm{aff}}=K_0$.
  \item The total initial slope is $K_p + K_H = K_p + Q_m K_0$;
        the universal workflow separates $K_p$ from $K_H$ by
        detecting the persistent linear trend at high~$C$.
\end{enumerate}
\end{proposition}

\begin{proof}
Parts (1)--(4) follow from Proposition~\ref{prop:langmuir}
applied to the saturating component.  Part (5) is the operational
strategy: at high $C$, $q\approx K_p C + Q_m$ (linear with
intercept $Q_m$ and slope $K_p$), so $K_p$ is identifiable from
the high-$C$ slope.  After subtracting $K_p C$, the residual
is a standard Langmuir.
\end{proof}

% ------ 3.7 Jovanovic ------
\subsection{Jovanovic isotherm}

\begin{proposition}[Jovanovic]
\label{prop:jovanovic}
The Jovanovic isotherm~\cite{Jovanovic1969}
\begin{equation}
  q(C) = q_{\max}\bigl(1 - e^{-bC}\bigr),
  \qquad b>0,\; q_{\max}>0,
  \label{eq:jovanovic}
\end{equation}
arises from the heterogeneous Langmuir integral with a specific
$\rho(K)$ when the local isotherm kernel is the Jovanovic kernel
$\varphi_J(C,K) = 1 - e^{-KC}$ rather than the Langmuir kernel.
Its invariants are:
\begin{enumerate}[nosep]
  \item $Q_{\max} = q_{\max}$ (finite for all $b$).
  \item $K_H = q_{\max}\,b$ (finite).
  \item $K_{\mathrm{aff}} = b$.
\end{enumerate}
\end{proposition}

\begin{proof}
\textit{Part (1):}
\[
  \lim_{C\to\infty} q_{\max}(1 - e^{-bC}) = q_{\max}\cdot(1 - 0) = q_{\max}.
\]

\textit{Part (2):}  The initial slope is
\[
  \frac{dq}{dC}\bigg|_{C=0}
  = q_{\max}\,b\,e^{-bC}\bigg|_{C=0}
  = q_{\max}\,b.
\]

\textit{Part (3):} $K_{\mathrm{aff}} = K_H / Q_{\max} = b$.
\end{proof}

\begin{remark}[Jovanovic kernel vs Langmuir kernel]
The Jovanovic kernel $\varphi_J(C,K) = 1 - e^{-KC}$ satisfies the
conditions (K1)--(K4) of Theorem~\ref{thm:kernel}: $\varphi_J(0,K)=0$,
$\partial\varphi_J/\partial C|_{C=0} = K$, and
$\lim_{C\to\infty}\varphi_J = 1$.  Therefore, the invariants
$Q_{\max}$, $K_H$, and $K_{\mathrm{aff}}$ are the same regardless of
whether the Langmuir or Jovanovic local mechanism is assumed.
The Jovanovic model describes a mobile adsorbed phase where molecules
undergo lateral diffusion, as opposed to the localized Langmuir model.
The site-energy distribution for the single-site Jovanovic model
corresponds to $\rho(K) = q_{\max}\,\delta(K - b)$ when the Jovanovic
kernel is used.

The energy distribution obtained by concentration normalisation
$C_e = C^* e^{-E/RT}$ is a Gumbel distribution (type-I extreme
value), with standard deviation
\begin{equation}
  \sigma_E^{\mathrm{(Jov)}} = \frac{\pi}{\sqrt{6}}\,RT,
  \label{eq:sigmaE_jov_v5}
\end{equation}
which is fixed by temperature alone and does not depend on the
parameter~$b$.  This contrasts with the Langmuir model
($\sigma_E = 0$, Dirac delta) and the Sips model
($\sigma_E = n\pi RT/\sqrt{3}$, adjustable).
\end{remark}

% ------ 3.8 UNILAN ------
\subsection{UNILAN isotherm}

\begin{proposition}[UNILAN]
\label{prop:unilan}
Let the affinity distribution be uniform in $\ln K$ over the interval
$[\ln b_{\min},\, \ln b_{\max}]$ with $0 < b_{\min} < b_{\max}$:
\begin{equation}
  f(K) = \frac{1}{K\,\ln(b_{\max}/b_{\min})},
  \qquad b_{\min} \le K \le b_{\max}.
  \label{eq:unilan_dist}
\end{equation}
Then the overall isotherm is
\begin{equation}
  q(C) = \frac{Q_{\max}}{2s}
  \ln\!\left(\frac{1 + b_{\max}\,C}{1 + b_{\min}\,C}\right),
  \qquad s = \frac{1}{2}\ln\frac{b_{\max}}{b_{\min}},
  \label{eq:unilan}
\end{equation}
and the invariants are:
\begin{enumerate}[nosep]
  \item $Q_{\max}$ is finite and correctly recovered.
  \item $K_H = \displaystyle
        \frac{Q_{\max}(b_{\max} - b_{\min})}
             {\ln(b_{\max}/b_{\min})}$ (finite).
  \item $K_{\mathrm{aff}} = \displaystyle
        \frac{b_{\max} - b_{\min}}{\ln(b_{\max}/b_{\min})}$
        \quad (logarithmic mean of $b_{\min}$ and $b_{\max}$).
\end{enumerate}
\end{proposition}

\begin{proof}
The derivation is identical to that of the Temkin isotherm
(Proposition~\ref{prop:temkin}) with
$K_{\min} \equiv b_{\min}$ and $K_{\max} \equiv b_{\max}$.
The uniform distribution in $\ln K$ gives
\[
  q(C) = Q_{\max}\int_{b_{\min}}^{b_{\max}}
  \frac{1}{K\ln(b_{\max}/b_{\min})}
  \cdot\frac{KC}{1+KC}\,dK
  = \frac{Q_{\max}}{\ln(b_{\max}/b_{\min})}
  \bigl[\ln(1+b_{\max}C) - \ln(1+b_{\min}C)\bigr],
\]
which is Eq.~\eqref{eq:unilan} with $2s = \ln(b_{\max}/b_{\min})$.
The saturation limit follows from
$\ln(1+b_{\max}C) - \ln(1+b_{\min}C) \to \ln(b_{\max}/b_{\min})$ as
$C\to\infty$.  The Henry constant and effective affinity are computed
exactly as in Proposition~\ref{prop:temkin}.
\end{proof}

\begin{remark}[Difference between UNILAN and Temkin]
The UNILAN and Temkin models share the same mathematical structure
(uniform energy distribution), but they are conventionally parameterised
differently.  The Temkin model is typically presented in its
intermediate-coverage approximation $q \approx (RT/b_T)\ln(a_T C)$,
whereas the UNILAN model retains the full logarithmic form.  In the
universal-route framework, both yield the same invariants, and the
effective affinity is the logarithmic mean of the boundary affinities:
\begin{equation}
  K_{\mathrm{aff}}^{\mathrm{(UNILAN)}}
  = \frac{b_{\max} - b_{\min}}{\ln(b_{\max}/b_{\min})}
  \equiv K_{\mathrm{aff}}^{\mathrm{(Temkin)}}.
\end{equation}
The energy-distribution width is
\begin{equation}
  \sigma_E^{\mathrm{(UNILAN)}}
  = \frac{RT}{2\sqrt{3}}\ln\frac{b_{\max}}{b_{\min}}
  = \frac{s\,RT}{\sqrt{3}}.
  \label{eq:sigmaE_unilan_v5}
\end{equation}
\end{remark}

% ------ 3.9 Dubinin-Radushkevich / Dubinin-Astakhov ------
\subsection{Dubinin--Radushkevich (D-R) and Dubinin--Astakhov (D-A)
isotherms}

\begin{proposition}[D-R / D-A --- partial compatibility]
\label{prop:DA}
The Dubinin--Astakhov isotherm
\begin{equation}
  q(C) = Q_{\max}\exp\!\left[
    -\left(\frac{\varepsilon}{E_a}\right)^{\!n}\right],
  \qquad \varepsilon = RT\ln\frac{C_s}{C},
  \label{eq:DA}
\end{equation}
where $C_s$ is the saturation concentration, $E_a$ is the
characteristic energy, and $n$ is the D-A exponent ($n=2$ for D-R),
has the following properties:
\begin{enumerate}[nosep]
  \item $Q_{\max}$ is finite: $\lim_{C\to C_s^-} q(C) = Q_{\max}$.
  \item $K_H = \lim_{C\to 0} dq/dC = +\infty$
        \quad (A3 violated for all $n \ge 1$).
  \item The model operates in the Polanyi potential framework, not in
        the Langmuir-kernel framework.  The ``affinity'' is encoded in
        $E_a$ rather than in a Langmuir constant.
\end{enumerate}
\end{proposition}

\begin{proof}
\textit{Part (1):}
As $C\to C_s$, $\varepsilon = RT\ln(C_s/C) \to 0$, so
$\exp[-(0/E_a)^n] = 1$ and $q\to Q_{\max}$.

\textit{Part (2):}
Write $q(C) = Q_{\max}\exp\![-A(\ln(C_s/C))^n]$ with $A = (RT/E_a)^n$.
Then
\[
  \frac{dq}{dC} = Q_{\max}\,A\,n\,\frac{[\ln(C_s/C)]^{n-1}}{C}
  \cdot\exp\![-A(\ln(C_s/C))^n].
\]
As $C\to 0^+$, $\ln(C_s/C)\to\infty$.  The exponential decay
$\sim\exp[-A(\ln(C_s/C))^n]$ competes with $1/C$.  For the D-R case
($n=2$), the exponential decays as $\exp[-A(\ln C_s/C)^2]$, which is
a very flat function of $1/C$.  One can verify that
\[
  \lim_{C\to 0^+}\frac{dq}{dC}
  = \lim_{u\to\infty} Q_{\max}An\,u^{n-1}e^u \cdot e^{-Au^n} = +\infty
  \qquad\text{for } n\ge 1,
\]
where $u = \ln(C_s/C)$ and $e^u = C_s/C$ grows faster than the
Gaussian-like suppression.

\textit{Part (3):}
The D-R/D-A model is derived from Polanyi's potential theory, where
the adsorption space is filled according to adsorption potential
$\varepsilon$, not from a superposition of Langmuir sites.
Therefore, the heterogeneous Langmuir integral representation
does not apply directly.  The universal route can still extract
an operational $\widehat{Q}_{\max}$ from the saturation plateau,
but $K_{\mathrm{aff}}$ will be range-dependent because the true
Henry constant is infinite.
\end{proof}

\begin{remark}[Operational use of D-R/D-A within the universal route]
Although the D-R/D-A model lies outside the strict heterogeneous
Langmuir framework, $Q_{\max}$ is well-defined and finite.
The mean adsorption energy $E_a$ plays the role of an affinity
descriptor within the Polanyi framework.  For cross-model comparison,
the energy-scale mapping
\begin{equation}
  E_{\mathrm{aff}}^{\mathrm{(DA)}}
  = E_a\,\Gamma\!\left(1 + \frac{1}{n}\right)
  \label{eq:Eaff_DA}
\end{equation}
(where $\Gamma$ is the gamma function) gives the mean adsorption
energy under the D-A distribution.  For D-R ($n=2$):
$E_{\mathrm{aff}}^{\mathrm{(DR)}} = E_a\sqrt{\pi}/2 \approx 0.886\,E_a$.
\end{remark}

% ------ 3.10 Hill ------
\subsection{Hill isotherm}

\begin{proposition}[Hill --- cooperative case]
\label{prop:hill}
The Hill isotherm
\begin{equation}
  q(C) = \frac{Q_{\max}\,(C/K_D)^{n_H}}{1 + (C/K_D)^{n_H}},
  \qquad n_H > 0,\; K_D > 0,
  \label{eq:hill}
\end{equation}
where $n_H$ is the Hill coefficient and $K_D$ is the half-saturation
concentration, has the following properties with respect to the
universal route:
\begin{enumerate}[nosep]
  \item $Q_{\max}$ is finite for all $n_H > 0$.
  \item For $n_H = 1$: reduces to Langmuir with $K_L = 1/K_D$,
        and all invariants are well-defined.
  \item For $0 < n_H < 1$ (anti-cooperative / heterogeneous):
        $K_H = +\infty$ (first moment diverges, same as Sips with
        $m = n_H$), but $Q_{\max}$ is finite.
        The isotherm is mathematically identical to the Sips model
        with $K_S = K_D^{-1}$ and $m = n_H$.
  \item For $n_H > 1$ (positive cooperativity):
        $K_H = 0$ (the initial slope vanishes), producing a sigmoidal
        isotherm.  This case violates Assumption A1 (independent sites)
        and is excluded from the universal route.
\end{enumerate}
\end{proposition}

\begin{proof}
\textit{Part (1):}
$(C/K_D)^{n_H}/[1+(C/K_D)^{n_H}] \to 1$ as $C\to\infty$ for any
$n_H > 0$.

\textit{Part (2):}
For $n_H = 1$: $q = Q_{\max}(C/K_D)/[1+C/K_D]
= Q_{\max}K_D^{-1}C/[1+K_D^{-1}C]$, which is the standard Langmuir
with $K_L = K_D^{-1}$.

\textit{Part (3):}
\[
  \frac{dq}{dC}\bigg|_{C=0}
  = \frac{Q_{\max}\,n_H}{K_D^{n_H}}\lim_{C\to 0^+} C^{n_H-1}
  = \begin{cases}
      +\infty & \text{if } n_H < 1, \\
      Q_{\max}/K_D & \text{if } n_H = 1, \\
      0 & \text{if } n_H > 1.
    \end{cases}
\]
For $n_H < 1$, the Hill isotherm is identical to the Sips isotherm
via $K_S = K_D^{-1}$ and $m = n_H$, as can be verified by direct
substitution into Eq.~\eqref{eq:sips_dist}.

\textit{Part (4):}
For $n_H > 1$, the vanishing initial slope produces an inflection
point in the isotherm (sigmoidal shape).  This implies cooperativity:
the adsorption of one molecule facilitates the adsorption of the next.
The isotherm then violates $d^2q/dC^2 < 0$ for all $C > 0$, which
is a necessary condition for heterogeneous Langmuir superposition.
\end{proof}

\begin{remark}[Effective affinity for Hill with $n_H < 1$]
For the anti-cooperative Hill isotherm ($n_H < 1$, identical to Sips),
an operational effective affinity can be defined as
\begin{equation}
  K_{\mathrm{aff}}^{\mathrm{eff}} = K_D^{-1/n_H} = K_S^{1/m},
\end{equation}
which reduces to $K_D^{-1} = K_L$ when $n_H = 1$ (Langmuir limit).
The corresponding energy-distribution width is
$\sigma_E = \pi RT / (\sqrt{3}\,n_H)$.
\end{remark}

% ======================================================================
\section{Universal low-concentration limit}
\label{sec:henry}
% ======================================================================

\begin{proposition}[Henry law as universal limit]
\label{prop:henry}
For any distribution $\rho(K)\ge 0$ satisfying A2 and A3 (with
$K_p\ge 0$):
\begin{equation}
  q(C) = (K_p + K_H)\,C + O(C^2) \qquad\text{as } C\to 0.
  \label{eq:henry_universal}
\end{equation}
\end{proposition}

\begin{proof}
The Langmuir kernel admits the Taylor expansion
$KC/(1+KC) = KC - K^2 C^2 + O(C^3)$.  Under assumption A3,
$\int_0^\infty K\rho(K)\,dK = K_H<\infty$, so by dominated
convergence we may interchange the integral and the limit:
\begin{align*}
  q(C) &= K_p C + \int_0^\infty \rho(K)\bigl[KC - K^2 C^2
          + O(C^3)\bigr]\,dK \\
       &= K_p C + K_H\,C - \left(\int_0^\infty K^2\rho(K)\,dK\right)C^2
          + O(C^3).
\end{align*}
If additionally the second moment
$\int K^2\rho(K)\,dK<\infty$, the $O(C^2)$ coefficient is
well-defined and equals $-Q_{\max}\langle K^2\rangle_f$.
\end{proof}

\begin{corollary}
The Henry constant $K_H$ is a model-independent invariant: it equals
$Q_{\max}\langle K\rangle_f$ regardless of the shape of $f(K)$,
as long as the first moment exists.
\end{corollary}

% ======================================================================
\section{Kernel-independence theorem}
\label{sec:kernel}
% ======================================================================

The invariants $Q_{\max}$ and $K_H$ (hence $K_{\mathrm{aff}}$) were
derived using the Langmuir kernel $\varphi(C,K)=KC/(1+KC)$.
The following theorem shows they are \emph{independent} of the
specific local isotherm mechanism.

\begin{theorem}[Kernel independence of asymptotic invariants]
\label{thm:kernel}
Let $\varphi(C,K)$ be \emph{any} local isotherm kernel satisfying:
\begin{enumerate}[nosep]
  \item[(K1)] $\varphi(0,K)=0$ for all $K>0$
        \quad (no uptake at zero concentration),
  \item[(K2)] $\displaystyle\frac{\partial\varphi}{\partial C}
        \bigg|_{C=0}=K$ for all $K>0$
        \quad (Henry-law onset with constant $K$),
  \item[(K3)] $\displaystyle\lim_{C\to\infty}\varphi(C,K)=1$
        for all $K>0$
        \quad (saturation of each site class),
  \item[(K4)] $0\le\varphi(C,K)\le 1$ for all $C,K\ge 0$
        \quad (bounded).
\end{enumerate}
Consider the generalised heterogeneous isotherm
\begin{equation}
  q(C)=\int_0^\infty\rho(K)\,\varphi(C,K)\,dK.
  \label{eq:general_kernel}
\end{equation}
Then, under assumptions A2 and A3:
\begin{equation}
  \lim_{C\to\infty}q(C)=Q_{\max},\qquad
  \frac{dq}{dC}\bigg|_{C=0}=K_H,\qquad
  K_{\mathrm{aff}}=\frac{K_H}{Q_{\max}}.
  \label{eq:kernel_indep}
\end{equation}
\end{theorem}

\begin{proof}
\textit{Saturation limit:}
By (K3), $\varphi(C,K)\to 1$ for each $K>0$ as $C\to\infty$.
By (K4), $|\rho(K)\varphi(C,K)|\le\rho(K)$, which is integrable by
A2.  By the dominated convergence theorem:
\[
  \lim_{C\to\infty}q(C) = \int_0^\infty \rho(K)\cdot 1\,dK
  = Q_{\max}.
\]

\textit{Henry slope:}
By (K1) and (K2), $\varphi(C,K)/C\to K$ as $C\to 0$.
By (K4), $\rho(K)\varphi(C,K)/C\le\rho(K)\cdot K$ (for $C$
small enough), which is integrable by A3.  By dominated convergence:
\[
  \lim_{C\to 0}\frac{q(C)}{C}
  = \int_0^\infty \rho(K)\cdot K\,dK = K_H.
\]

\textit{Effective affinity:}
$K_{\mathrm{aff}}=K_H/Q_{\max}$ by definition.
\end{proof}

\begin{remark}[Examples of valid kernels]
In addition to the Langmuir kernel $\varphi_L=KC/(1+KC)$,
Theorem~\ref{thm:kernel} applies to:
\begin{itemize}[nosep]
  \item \textbf{Jovanovic kernel:}
        $\varphi_J(C,K)=1-e^{-KC}$.
        Check: $\varphi_J(0,K)=0$, $\varphi_J'(0)=K$,
        $\varphi_J(\infty)=1$.
  \item \textbf{Volmer kernel:}
        Any kernel derived from a 2-D equation of state for the
        adsorbed phase, provided it satisfies (K1)--(K4).
\end{itemize}
The NNLS distribution recovery \emph{does} depend on the kernel
choice, but the asymptotic invariants $Q_{\max}$, $K_H$, and
$K_{\mathrm{aff}}$ do not.
\end{remark}

% ======================================================================
\section{Degenerate cases: divergent invariants}
\label{sec:degenerate}
% ======================================================================

% ------ 6.1 Freundlich ------
\subsection{Pure Freundlich isotherm}

\begin{proposition}[Freundlich --- doubly degenerate]
\label{prop:freundlich}
The Freundlich isotherm $q(C)=K_F C^{n}$ with $0<n<1$ cannot be
represented by the heterogeneous Langmuir integral with finite
$Q_{\max}$ and $K_H$.  Specifically:
\begin{enumerate}[nosep]
  \item $Q_{\max}=\lim_{C\to\infty}q(C)=+\infty$
        \quad (A2 violated).
  \item $K_H=\lim_{C\to 0}dq/dC=+\infty$
        \quad (A3 violated).
  \item The underlying distribution $f(K)\propto K^{n-2}$ has both
        zeroth and first moments infinite.
\end{enumerate}
\end{proposition}

\begin{proof}
\textit{Part (1):} $\lim_{C\to\infty}K_F C^n=+\infty$ for $n>0$.

\textit{Part (2):}
$dq/dC=K_F n C^{n-1}\to+\infty$ as $C\to 0^+$ for $n<1$.

\textit{Part (3):}
The Freundlich isotherm is the low-coverage limit of the Sips
isotherm with $m=n$ and $Q_{\max}\to\infty$ (see
Proposition~\ref{prop:sips}).
The Sips distribution $f(K)$ has a tail $\sim K^{-m-1}$,
whose zeroth moment over $[0,\infty)$ diverges when $Q_{\max}$ is
not finite.
\end{proof}

\begin{remark}[Operational treatment of Freundlich data]
When the universal workflow is applied to data that follow a
Freundlich isotherm over a finite range $[C_{\min},C_{\max}]$:
\begin{itemize}[nosep]
  \item $\widehat{Q}_{\max}$ will be approximately $K_F C_{\max}^n$
        (the value at the highest measured concentration).
  \item $\widehat{K}_H$ will depend strongly on $C_{\min}$.
  \item The sensitivity diagnostic will flag
        $K_{\mathrm{aff}}$ as range-dependent.
  \item These operational estimates are \emph{not} material constants
        but depend on the experimental window.
\end{itemize}
\end{remark}

% ------ 6.2 Redlich-Peterson ------
\subsection{Redlich--Peterson isotherm}

\begin{proposition}[Redlich--Peterson]
\label{prop:rp}
The Redlich--Peterson isotherm~\cite{Redlich1959}
\begin{equation}
  q(C) = \frac{K_R\,C}{1+a_R\,C^{\beta}},
  \qquad 0<\beta\le 1,
  \label{eq:rp}
\end{equation}
has the following properties:
\begin{enumerate}[nosep]
  \item For $\beta=1$: reduces to Langmuir with
        $Q_{\max}=K_R/a_R$ and $K=a_R$.
  \item For $\beta<1$:
        $Q_{\max}=\lim_{C\to\infty}q(C)=+\infty$ (A2 violated),
        but $K_H=K_R$ is finite.
  \item The Redlich--Peterson isotherm with $\beta<1$ does not arise
        from the heterogeneous Langmuir integral with finite
        $Q_{\max}$.
\end{enumerate}
\end{proposition}

\begin{proof}
\textit{Part (1):} Set $\beta=1$:
$q=K_R C/(1+a_R C)=(K_R/a_R)\cdot a_R C/(1+a_R C)$, which is
Langmuir.

\textit{Part (2):}
$\lim_{C\to\infty}K_R C/(a_R C^\beta)
=(K_R/a_R)C^{1-\beta}\to\infty$ for $\beta<1$.
At $C\to 0$: $q\approx K_R C/(1+0)=K_R C$, so $K_H=K_R$.

\textit{Part (3):}
Since $Q_{\max}=\infty$, the isotherm cannot be written as
Eq.~\eqref{eq:master} with $\int\rho(K)\,dK<\infty$.
\end{proof}

% ======================================================================
\section{Models excluded by assumption A1}
\label{sec:excluded}
% ======================================================================

The following models violate Assumption A1
(independent-site superposition) and are therefore \emph{outside}
the scope of the universal route:

\begin{enumerate}[nosep]
  \item \textbf{BET isotherm}~\cite{BET1938}:
        Multilayer adsorption.  Each adsorbed molecule provides a
        new site for the next layer, violating independence.
        $Q_{\max}$ in the BET sense refers to monolayer coverage,
        but the isotherm diverges as $C\to C_s$ (saturation
        concentration).

  \item \textbf{Hill isotherm}
        $q=Q_{\max}(KC)^{n_H}/(1+(KC)^{n_H})$ with $n_H>1$:
        Positive cooperativity produces a sigmoidal isotherm.
        At $C=0$, $dq/dC=0$ for $n_H>1$, meaning $K_H=0$.
        The inflection point violates the monotone-concave shape
        expected from heterogeneous Langmuir superposition.

  \item \textbf{Fowler--Guggenheim isotherm}:
        Lateral interactions between adsorbed molecules
        ($q$-dependent free energy) produce a non-Langmuir local
        isotherm.  For strong enough interactions, the isotherm
        can become S-shaped or exhibit hysteresis.
\end{enumerate}

\textbf{Diagnostic:}
If the data show a sigmoidal shape (inflection point), the isotherm
lies outside the heterogeneous Langmuir class, and the universal
route should not be applied.  A simple test: verify that $d^2 q/dC^2<0$
for all measured $C>0$.

% ======================================================================
\section{Summary: classification of isotherm models}
\label{sec:summary}
% ======================================================================

Table~\ref{tab:summary} classifies the main adsorption isotherm
models with respect to the universal route.

\begin{table}[htbp]
\centering
\caption{Convergence of the universal route to classical isotherms.
``$\checkmark$'' = finite and correctly recovered;
``$\infty$'' = divergent (assumption violated);
``---'' = not applicable.}
\label{tab:summary}
\small
\begin{tabular}{@{}lcccccl@{}}
\toprule
\textbf{Model} & \textbf{A1--A4} & $Q_{\max}$ & $K_H$
  & $K_{\mathrm{aff}}$ & $K_p$ & \textbf{Proposition} \\
\midrule
Langmuir
  & all & $Q_m$ & $Q_m K$ & $K$
  & $0$ & \ref{prop:langmuir} \\[2pt]
Multi-Langmuir ($J$-site)
  & all & $\sum Q_j$ & $\sum Q_j K_j$
  & $\frac{\sum Q_j K_j}{\sum Q_j}$
  & $0$ & \ref{prop:multilang} \\[2pt]
Sips ($m=1$)
  & all & $Q_m$ & $Q_m K_s$ & $K_s$
  & $0$ & \ref{prop:sips} \\[2pt]
Sips ($m<1$)
  & A3 viol. & $Q_m\;\checkmark$ & $\infty$
  & oper. & $0$ & \ref{prop:sips} \\[2pt]
T\'oth ($0<t\le 1$)
  & all & $Q_m$ & $Q_m K_T$ & $K_T$
  & $0$ & \ref{prop:toth} \\[2pt]
Temkin
  & all & $Q_m$ & $Q_m K_{\log}$
  & $K_{\log}$ & $0$ & \ref{prop:temkin} \\[2pt]
Dual-mode
  & all & $Q_m$ & $Q_m K$ & $K$
  & $K_p$ & \ref{prop:dualmode} \\[2pt]
Jovanovic
  & all & $q_{\max}$ & $q_{\max}b$ & $b$
  & $0$ & \ref{prop:jovanovic} \\[2pt]
UNILAN
  & all & $Q_m$ & $Q_m K_{\log}$
  & $K_{\log}$ & $0$ & \ref{prop:unilan} \\[2pt]
D-R/D-A
  & A3 viol. & $Q_m\;\checkmark$ & $\infty$
  & --- & $0$ & \ref{prop:DA} \\[2pt]
Hill ($n_H<1$)
  & A3 viol. & $Q_m\;\checkmark$ & $\infty$
  & oper. & $0$ & \ref{prop:hill} \\[2pt]
Hill ($n_H=1$)
  & all & $Q_m$ & $Q_m/K_D$ & $1/K_D$
  & $0$ & \ref{prop:hill} \\[2pt]
Freundlich
  & A2,A3 viol. & $\infty$ & $\infty$
  & --- & --- & \ref{prop:freundlich} \\[2pt]
Redlich--Peterson ($\beta<1$)
  & A2 viol. & $\infty$ & $K_R\;\checkmark$
  & --- & $0$ & \ref{prop:rp} \\[2pt]
\midrule
BET & A1 viol. & --- & --- & --- & ---
  & Sect.~\ref{sec:excluded} \\
Hill ($n_H>1$)
  & A1 viol. & $Q_m$ & $0$ & $0$
  & --- & Sect.~\ref{sec:excluded} \\
Fowler--Guggenheim
  & A1 viol. & --- & --- & --- & ---
  & Sect.~\ref{sec:excluded} \\
\bottomrule
\end{tabular}

\medskip
$K_{\log}\equiv(K_{\max}-K_{\min})/\ln(K_{\max}/K_{\min})$
\quad (logarithmic mean).
\end{table}

\textbf{Key conclusions:}
\begin{enumerate}[nosep]
  \item Every isotherm within the heterogeneous Langmuir class
        (A1 satisfied) with finite $Q_{\max}$ and $K_H$
        (A2--A3 satisfied) yields exact convergence of the
        universal invariants to the model parameters.
  \item The T\'oth isotherm is the broadest single-equation model
        that satisfies all assumptions: it accommodates asymmetric
        heterogeneity while keeping $K_{\mathrm{aff}}=K_T$
        well-defined for all $t\in(0,1]$.
  \item The Sips isotherm with $m<1$ produces a finite
        $Q_{\max}$ but an infinite $K_H$; the universal workflow
        correctly recovers $Q_{\max}$ and flags the
        $K_{\mathrm{aff}}$ estimate as range-dependent.
  \item The Freundlich and Redlich--Peterson ($\beta<1$) isotherms
        have infinite $Q_{\max}$; the universal workflow gives
        operational estimates that are inherently range-dependent.
  \item The Jovanovic model ($K_{\mathrm{aff}} = b$) and UNILAN model
        ($K_{\mathrm{aff}} = K_{\log}$) satisfy all assumptions
        and converge exactly, the former using a different local
        kernel (Theorem~\ref{thm:kernel} guarantees invariance).
  \item The D-R/D-A model has finite $Q_{\max}$ but infinite $K_H$
        (Polanyi framework); the affinity energy $E_a$ serves as an
        alternative affinity descriptor.
  \item The Hill isotherm with $n_H < 1$ is equivalent to Sips;
        with $n_H > 1$ it exhibits cooperativity and is excluded.
  \item Models violating A1 (BET, Hill with $n_H>1$,
        Fowler--Guggenheim) produce isotherms outside the
        concave-Langmuir class and are excluded.
\end{enumerate}

% ======================================================================
\section{Numerical verification}
\label{sec:numerical}
% ======================================================================

The companion MATLAB script
\texttt{universal\_route\_v4\_limits.m} verifies the
convergence results numerically.  For each model in
Table~\ref{tab:summary} (excluding the excluded models), synthetic
data with 5\% relative Gaussian noise are generated, the universal
workflow is applied, and the estimated invariants are compared with
the analytical values.

% ------ Figure 1 ------
\subsection{Isotherm reconstruction for all models}

Figure~\ref{fig:v4models} shows the reconstruction for seven
representative models.  In each panel, the grey curve is the true
generating model, black circles are the noisy data, the red dashed
curve is the effective Langmuir reconstruction using
$\widehat{Q}_{\max}$ and $\widehat{K}_{\mathrm{aff}}$, and the
green horizontal line marks $\widehat{Q}_{\max}$.

For the Langmuir, T\'oth, Temkin, and dual-mode models, the
reconstruction closely follows the true curve and the estimated
invariants match the analytical values within the noise level.
For the Sips model, $Q_{\max}$ is well estimated but
$K_{\mathrm{aff}}$ is operational (finite but range-dependent).
For the Freundlich model, both invariants are operational
approximations to the (infinite) true values.

\begin{figure}[htbp]
  \centering
  \includegraphics[width=\textwidth]{fig_v4_models.pdf}
  \caption{Isotherm reconstruction for seven models using the
  universal route.  Grey: true model; black circles: noisy data;
  red dashed: effective Langmuir reconstruction; green: estimated
  $\widehat{Q}_{\max}$.  The bottom-right panel shows the relative
  errors $|\widehat{Q}_{\max}-Q_{\max}^{\mathrm{true}}|
  /Q_{\max}^{\mathrm{true}}$ (blue) and
  $|\widehat{K}_{\mathrm{aff}}-K_{\mathrm{aff}}^{\mathrm{true}}|
  /K_{\mathrm{aff}}^{\mathrm{true}}$ (orange) for the models
  with finite invariants.}
  \label{fig:v4models}
\end{figure}

% ------ Figure 2 ------
\subsection{Convergence with decreasing noise}

Figure~\ref{fig:v4convergence} demonstrates that the estimation
errors in $Q_{\max}$ and $K_{\mathrm{aff}}$ decrease systematically
as the noise level is reduced, confirming consistency of the
universal workflow.

For models with well-defined invariants (Langmuir, T\'oth, Temkin,
dual-mode, double Langmuir), the relative error decreases roughly
linearly with the noise level on a log-log scale, indicating
$O(\sigma)$ convergence rate.
For the Sips model, $Q_{\max}$ converges but $K_{\mathrm{aff}}$
remains sensitive to the concentration range even at low noise,
confirming the first-moment divergence.

\begin{figure}[htbp]
  \centering
  \includegraphics[width=0.85\textwidth]{fig_v4_convergence.pdf}
  \caption{Relative error in $\widehat{Q}_{\max}$ (left) and
  $\widehat{K}_{\mathrm{aff}}$ (right) versus noise level
  for five models with finite invariants.  Each point is the
  median over 20 Monte Carlo realisations; error bars show
  the interquartile range.}
  \label{fig:v4convergence}
\end{figure}

% ======================================================================
\section{Implementation: fitting algorithm and code explanation}
\label{sec:implementation}
% ======================================================================

This section describes the computational implementation of the
universal route.  The algorithm is provided in both MATLAB
(\texttt{universal\_route\_corrected.m}) and Python
(\texttt{universal\_route\_corrected.py}).  We explain each step of
the workflow, the mathematical rationale behind the statistical
tests used, and the key design decisions.

% ------ 8.1 Overview ------
\subsection{Algorithm overview}

The fitting procedure takes as input:
\begin{itemize}[nosep]
  \item $\{C_{e,i}\}_{i=1}^{N}$: measured equilibrium concentrations.
  \item $\{q_{e,i}\}_{i=1}^{N}$: measured equilibrium uptakes.
  \item $\{\sigma_i\}_{i=1}^{N}$: uncertainties (standard deviations)
        on each $q_{e,i}$.
\end{itemize}
The output is the set of universal invariants
$\{\widehat{Q}_{\max},\, \widehat{K}_H,\,
\widehat{K}_{\mathrm{aff}},\, \widehat{K}_p\}$, together with
diagnostic information (Henry region size, plateau region size,
partition detection flag, $K_H$ sensitivity).

The algorithm proceeds in three main steps:
\begin{enumerate}
  \item \textbf{Partition test:} Detect whether a non-saturating
        linear component $K_p C$ is present (dual-mode behaviour).
  \item \textbf{Estimate $Q_{\max}$:} Identify the plateau region
        and extract the saturation capacity.
  \item \textbf{Estimate $K_H$:} Identify the Henry region and
        extract the initial slope.
\end{enumerate}

% ------ 8.2 Step 1: Partition test ------
\subsection{Step 1: Partition detection}

The partition test determines whether the high-concentration behaviour
is consistent with a saturating isotherm ($q \to Q_{\max}$) or with
a persistent linear trend ($q \approx K_p C + Q_{\max}$).

Three independent tests are applied:

\subsubsection{Test A: OLS $t$-test on the high-$C$ slope}

Select the $m$ highest-concentration points ($m = \min(12, \max(5,
\lfloor N/3\rfloor))$) and fit a linear model
\begin{equation}
  q_i = \beta_0 + \beta_1 C_{e,i} + \varepsilon_i,
  \qquad i \in \mathcal{I}_{\infty}.
  \label{eq:partition_linear}
\end{equation}
The fit is performed by \emph{ordinary} least squares (OLS),
not weighted least squares (WLS).  This is a deliberate choice:
when $\sigma_i \propto q_i$ (relative errors), WLS assigns the
smallest weights to the highest-$C$ points, which are precisely
where the partition signal is strongest.

The $t$-statistic for the slope is
\begin{equation}
  t_{\beta_1} = \frac{\hat{\beta}_1}{\mathrm{SE}(\hat{\beta}_1)},
  \qquad
  \mathrm{SE}(\hat{\beta}_1) = \sqrt{s^2 (X^T X)^{-1}_{22}},
  \label{eq:t_test_slope}
\end{equation}
where $s^2 = \sum r_i^2 / (m-2)$ is the residual variance.
The partition is detected if $t_{\beta_1} > t_{0.975}(m-2)$ and
$\hat{\beta}_1 > 0$.

\subsubsection{Test B: Ratio check}

Compare the mean of $q/C$ at high concentrations with the mean at
low concentrations:
\begin{equation}
  R_{\mathrm{high}} = \frac{1}{m}\sum_{i\in\mathcal{I}_\infty}
  \frac{q_i}{C_{e,i}},
  \qquad
  R_{\mathrm{low}} = \frac{1}{3}\sum_{i=1}^{3}
  \frac{q_i}{C_{e,i}}.
  \label{eq:ratio_check}
\end{equation}
For a pure saturating isotherm, $R_{\mathrm{high}} \ll R_{\mathrm{low}}$
(since $q/C \to Q_{\max}/C \to 0$).  The partition is flagged if
$R_{\mathrm{high}} > 0.10\,R_{\mathrm{low}}$.

\subsubsection{Test C: Derivative flatness}

Compute the numerical derivative $dq/dC$ from successive data pairs
at high~$C$.  For a partition model, $dq/dC \to K_p > 0$ (constant).
For a power-law (Freundlich), $dq/dC$ is monotonically decreasing.
The test checks:
\begin{enumerate}[nosep]
  \item The mean derivative exceeds a threshold:
        $\overline{dq/dC} > 0.02\,q_{\mathrm{mid}}/C_{\mathrm{mid}}$.
  \item The derivative is approximately flat: the ratio of median
        derivatives in the second and first halves of the high-$C$
        window satisfies $r > 0.70$.
\end{enumerate}

\subsubsection{Combined decision}

The partition is declared present if \emph{any} of the three tests
is positive:
\begin{equation}
  \texttt{use\_partition}
  = \texttt{Test A} \lor \texttt{Test B} \lor \texttt{Test C}.
\end{equation}

% ------ 8.3 Step 2: Qmax ------
\subsection{Step 2: Estimation of $Q_{\max}$}

\subsubsection{Case 1: No partition ($K_p = 0$)}

The plateau region $\mathcal{I}_\infty$ is identified by progressively
including high-$C$ points and testing whether the data are consistent
with a constant.  Starting from the last 3 points, we expand the
window and fit
\begin{equation}
  q_i = \alpha_0 + \alpha_1 \ln C_{e,i} + \varepsilon_i,
  \qquad i \in \mathcal{I}_\infty.
\end{equation}
A weighted $t$-test (with weights $w_i = 1/\sigma_i^2$) on the
coefficient $\alpha_1$ determines whether a logarithmic trend is
significant.  The plateau region is the largest window for which
$|t_{\alpha_1}| < t_{0.975}(\mathrm{df})$.

Once the plateau region is identified:
\begin{itemize}[nosep]
  \item If $n_{\mathrm{plateau}} \ge 3$: compute the weighted mean
        \begin{equation}
          \widehat{Q}_{\max}
          = \frac{\sum_{i\in\mathcal{I}_\infty} w_i\, q_i}
                 {\sum_{i\in\mathcal{I}_\infty} w_i}.
          \label{eq:Qmax_wmean}
        \end{equation}
  \item If $n_{\mathrm{plateau}} < 3$ (no clear plateau): extrapolate
        using a $1/C$ model,
        $q = \beta_0 + \beta_1/C$, and take
        $\widehat{Q}_{\max} = \hat{\beta}_0$.
\end{itemize}

\subsubsection{Case 2: Partition detected ($K_p > 0$)}

An iterative scheme jointly estimates $K_p$ and $Q_{\max}$:
\begin{enumerate}[nosep]
  \item \textbf{Initialise:} Fit $q = \beta_0 + \beta_1 C$ to the
        high-$C$ points.  Set $K_p^{(0)} = \max(\hat{\beta}_1, 0)$
        and $Q_{\max}^{(0)} = \max(\hat{\beta}_0, 0)$.
  \item \textbf{Iterate} (up to 30 times):
    \begin{enumerate}[nosep]
      \item Compute the saturating component:
            $q_{\mathrm{sat}} = q - K_p^{(k)} C$.
      \item Fit $q_{\mathrm{sat}} = \gamma_0 + \gamma_1/C$ to the
            high-$C$ points; update
            $Q_{\max}^{(k+1)} = \max(\hat{\gamma}_0, 0)$.
      \item Update $K_p$: at high $C$,
            $q - Q_{\max}^{(k+1)} \approx K_p C$, so
            \begin{equation}
              K_p^{(k+1)} = \max\!\left(
              \frac{\sum w_i C_i (q_i - Q_{\max}^{(k+1)})}
                   {\sum w_i C_i^2},\; 0\right).
            \end{equation}
    \end{enumerate}
  \item \textbf{Convergence:} Stop when
        $\max(|\Delta K_p|/K_p,\, |\Delta Q_{\max}|/Q_{\max}) < 10^{-4}$.
\end{enumerate}

% ------ 8.4 Step 3: KH ------
\subsection{Step 3: Estimation of $K_H$}

The Henry region $\mathcal{I}_0$ is identified by progressively
including low-$C$ points and testing for departure from linearity.
Starting from the first 3 points, we fit the linear model
\begin{equation}
  q_i = K_H\, C_{e,i},
  \qquad i \in \mathcal{I}_0
  \label{eq:henry_fit}
\end{equation}
(no intercept, since $q(0) = 0$), and the quadratic model
\begin{equation}
  q_i = \beta_1 C_{e,i} + \beta_2 C_{e,i}^2.
\end{equation}
An $F$-test compares the two:
\begin{equation}
  F = \frac{(\mathrm{SSE}_{\mathrm{lin}} - \mathrm{SSE}_{\mathrm{quad}})/1}
           {\mathrm{SSE}_{\mathrm{quad}} / (m-2)}.
  \label{eq:F_test}
\end{equation}
The Henry region is the largest window for which
$F < F_{0.95}(1, m-2)$, i.e., the quadratic term does not
significantly improve the fit.

Once the Henry region is identified, the Henry constant is estimated by
weighted regression through the origin:
\begin{equation}
  \widehat{K}_H
  = \frac{\sum_{i\in\mathcal{I}_0} w_i\, C_{e,i}\, q_i}
         {\sum_{i\in\mathcal{I}_0} w_i\, C_{e,i}^2}.
  \label{eq:KH_estimate}
\end{equation}

If a partition component was detected, $K_p C$ is subtracted before
computing $K_H$ (i.e., $q_{\mathrm{for\_henry}} = q - K_p C$).

\subsection{Step 4: Effective affinity and diagnostics}

The effective affinity is
\begin{equation}
  \widehat{K}_{\mathrm{aff}} = \frac{\widehat{K}_H}{\widehat{Q}_{\max}}.
\end{equation}

A convergence diagnostic compares $K_H$ estimated from the first
3~points with $K_H$ estimated from a slightly extended window
($n_{\mathrm{henry}} + 3$ points):
\begin{equation}
  \Delta K_H^{\mathrm{rel}}
  = \frac{|\widehat{K}_H^{(n+3)} - \widehat{K}_H^{(3)}|}
         {\widehat{K}_H^{(3)}}.
\end{equation}
If $\Delta K_H^{\mathrm{rel}} > 0.30$, the first moment of $f(K)$
may be divergent (Sips, Freundlich), and
$\widehat{K}_{\mathrm{aff}}$ should be reported as operational
(range-dependent).

% ------ 8.5 Bootstrap ------
\subsection{Bootstrap confidence intervals}

To quantify the uncertainty in $\widehat{Q}_{\max}$ and
$\widehat{K}_{\mathrm{aff}}$, a nonparametric bootstrap is used:
\begin{enumerate}[nosep]
  \item Draw $B$ bootstrap samples (typically $B = 1000$) by
        resampling $N$ data points $(C_{e,i}, q_{e,i}, \sigma_i)$
        with replacement.
  \item Apply the universal route to each bootstrap sample, obtaining
        $\{\widehat{Q}_{\max}^{(b)},\, \widehat{K}_{\mathrm{aff}}^{(b)}\}_{b=1}^{B}$.
  \item Compute the 95\% confidence intervals from the 2.5th and
        97.5th percentiles of the bootstrap distribution:
        \begin{equation}
          \mathrm{CI}_{95\%}(Q_{\max})
          = \bigl[\widehat{Q}_{\max}^{(0.025)},\;
                   \widehat{Q}_{\max}^{(0.975)}\bigr].
        \end{equation}
\end{enumerate}

% ------ 8.6 NNLS ------
\subsection{Non-parametric distribution recovery: Tikhonov-regularised NNLS}

As an independent check, the affinity distribution $\rho(K)$ can be
recovered non-parametrically from the data.  The continuous integral
Eq.~\eqref{eq:master} is discretised on $M$ logarithmically spaced
grid points $\{K_j\}_{j=1}^{M}$:
\begin{equation}
  q_i \approx K_p\,C_{e,i}
  + \sum_{j=1}^{M} q_{m,j}\,\frac{K_j C_{e,i}}{1 + K_j C_{e,i}},
  \qquad q_{m,j} \ge 0,
  \label{eq:NNLS_discrete}
\end{equation}
where $q_{m,j}$ represents the adsorption capacity associated with
affinity $K_j$.  In matrix form: $\mathbf{q} = \mathbf{A}\,\mathbf{q}_m$,
where $A_{ij} = K_j C_{e,i}/(1 + K_j C_{e,i})$.

The non-negativity constraint $q_{m,j} \ge 0$ is enforced by
non-negative least squares (NNLS).  To regularise the solution
(suppress oscillations), a Tikhonov penalty with a second-difference
operator $\mathbf{L}$ is added:
\begin{equation}
  \min_{\mathbf{q}_m \ge 0}\;
  \|\mathbf{W}^{1/2}(\mathbf{A}\mathbf{q}_m - \mathbf{q})\|^2
  + \lambda\,\|\mathbf{L}\mathbf{q}_m\|^2,
  \label{eq:NNLS_tikhonov}
\end{equation}
where $\mathbf{W} = \mathrm{diag}(1/\sigma_i^2)$ and $\lambda > 0$
is the regularisation parameter.

The optimal $\lambda$ is selected by the L-curve
method~\cite{Cernik1995}: the point of maximum curvature in the
$\log\|\mathbf{r}\|$ vs.\ $\log\|\mathbf{L}\mathbf{q}_m\|$ plot,
where $\mathbf{r} = \mathbf{W}^{1/2}(\mathbf{A}\mathbf{q}_m
- \mathbf{q})$.

The recovered distribution yields
\begin{equation}
  \widehat{Q}_{\max}^{\mathrm{(NNLS)}} = \sum_{j=1}^{M} q_{m,j},
  \qquad
  \widehat{K}_H^{\mathrm{(NNLS)}} = \sum_{j=1}^{M} K_j\, q_{m,j},
  \qquad
  \widehat{K}_{\mathrm{aff}}^{\mathrm{(NNLS)}}
  = \frac{\widehat{K}_H^{\mathrm{(NNLS)}}}
         {\widehat{Q}_{\max}^{\mathrm{(NNLS)}}}.
\end{equation}

% ------ 8.7 Code listing ------
\subsection{Annotated code: \texttt{universal\_route\_corrected} (MATLAB)}

The following listing shows the core function with annotations
linking each code block to the mathematical formulation above.

\begin{lstlisting}[caption={Core function of the universal route
  (MATLAB).  Line numbers refer to the mathematical derivations
  in Sections~\ref{sec:implementation}.1--\ref{sec:implementation}.4.},
  label=lst:core]
function res = universal_route_corrected(Ce, q, sigma)
    % Vectorise and sort by increasing concentration
    Ce = Ce(:); q = q(:); sigma = sigma(:);
    w = 1 ./ sigma.^2;           % inverse-variance weights
    N = length(Ce);
    [Ce, idx] = sort(Ce);
    q = q(idx); w = w(idx); sigma = sigma(idx);

    % ---- STEP 1: Partition test (Sec. 8.2) ----
    m_high = min(12, max(5, floor(N/3)));
    I_high = (N-m_high+1):N;

    % Test A: OLS t-test on slope [Eq. (8.2)]
    X_h = [ones(m_high,1), Ce(I_high)];
    beta_h = pinv(X_h'*X_h) * (X_h' * q(I_high));
    resid_h = q(I_high) - X_h * beta_h;
    s2_h = sum(resid_h.^2) / max(m_high-2, 1);
    se_slope = sqrt(max(s2_h*pinv(X_h'*X_h), 0));
    se_slope = se_slope(2,2);
    t_slope = beta_h(2) / se_slope;
    % Compare with t_critical at 97.5% (two-sided)
    use_partition_ttest = (t_slope > t_crit) && (beta_h(2) > 0);

    % Test B: ratio check [Eq. (8.3)]
    ratio_high = q(I_high) ./ Ce(I_high);
    ratio_low  = q(1:3) ./ Ce(1:3);
    use_partition_ratio = mean(ratio_high) > 0.10*mean(ratio_low);

    % Test C: derivative flatness (Sec. 8.2.3)
    % ... (flatness check on numerical dq/dC)

    use_partition = test_A || test_B || test_C;

    % ---- STEP 2: Estimate Qmax (Sec. 8.3) ----
    if ~use_partition
        % Expand plateau window, t-test on log(C) trend
        % Weighted mean of plateau points [Eq. (8.5)]
        Qmax = sum(w(I_plat).*q(I_plat)) / sum(w(I_plat));
        Kp = 0;
    else
        % Iterative Kp/Qmax decomposition (Sec. 8.3.2)
        % ... 30 iterations of alternating updates
    end

    % ---- STEP 3: Estimate KH via F-test (Sec. 8.4) ----
    % Expand Henry window, F-test: linear vs quadratic
    % [Eq. (8.7)]
    for m_try = 3:floor(N/2)
        % Fit q = KH*C (linear, no intercept)
        % Fit q = b1*C + b2*C^2 (quadratic)
        % F = (SSE_lin - SSE_quad)/1 / (SSE_quad/(m-2))
        % Stop when F > F_critical
    end
    KH = sum(w.*Ce.*q) / sum(w.*Ce.^2);  % [Eq. (8.8)]

    % ---- STEP 4: Affinity and diagnostics ----
    Kaff = KH / Qmax;
    % Convergence diagnostic: compare KH from 3 pts vs n+3 pts
    res = struct('Qmax',Qmax, 'KH',KH, 'Kaff',Kaff, 'Kp',Kp);
end
\end{lstlisting}

\subsection{Python implementation}

A fully equivalent Python implementation is provided in
\texttt{universal\_route\_corrected.py}, using \texttt{numpy} for
linear algebra, \texttt{scipy.optimize.nnls} for the non-negative
least squares solver, and \texttt{scipy.stats} for the $t$ and $F$
distribution critical values (replacing the look-up tables used in
the MATLAB version).

The key differences from the MATLAB version are:
\begin{itemize}[nosep]
  \item \texttt{numpy.linalg.pinv} replaces MATLAB's \texttt{pinv}
        for pseudoinverse computation.
  \item \texttt{scipy.stats.t.ppf(0.975, df)} replaces the
        hard-coded $t$-critical value table.
  \item \texttt{scipy.stats.f.ppf(0.95, 1, df2)} replaces the
        hard-coded $F$-critical value table.
  \item Bootstrap and NNLS functions are modularised as separate
        functions with the same interface.
\end{itemize}

Both implementations have been verified to produce identical results
(within floating-point precision) on the same synthetic datasets, as
shown in Section~\ref{sec:numerical}.

% ------ 8.8 Statistical helpers ------
\subsection{Statistical helper functions}

The MATLAB implementation avoids dependence on the Statistics Toolbox
by using hard-coded look-up tables for critical values:

\subsubsection{$t$-distribution critical values}

The function \texttt{t\_critical\_975(df)} returns $t_{0.975}(\nu)$
for degrees of freedom $\nu$ by linear interpolation of a table
covering $\nu \in \{1, 2, \ldots, 10, 12, 15, 20, 25, 30, 40, 60,
120, \infty\}$.  The asymptotic value is $t_{0.975}(\infty) = 1.960$.

\subsubsection{$F$-distribution critical values}

The function \texttt{f\_critical\_95(1, df2)} returns
$F_{0.95}(1, \nu_2)$ for numerator $\mathrm{df}_1 = 1$ (since the
$F$-test always compares a linear model with one extra parameter).
The asymptotic value is $F_{0.95}(1, \infty) = 3.841 = \chi^2_{0.95}(1)$.

% ======================================================================
\section*{References}
% ======================================================================

\begin{thebibliography}{99}

\bibitem{Sips1948}
R.~Sips, ``On the structure of a catalyst surface,''
\textit{J.\ Chem.\ Phys.}\ \textbf{16}, 490--495 (1948).

\bibitem{Jaroniec1988}
M.~Jaroniec and R.~Madey,
\textit{Physical Adsorption on Heterogeneous Solids},
Elsevier, Amsterdam (1988).

\bibitem{Rudzinski1992}
W.~Rudzinski and D.H.~Everett,
\textit{Adsorption of Gases on Heterogeneous Surfaces},
Academic Press, London (1992).

\bibitem{Klotz1971}
I.M.~Klotz and D.L.~Hunston,
``Properties of graphical representations of multiple classes of
binding sites,''
\textit{Biochemistry} \textbf{10}, 3065--3069 (1971).

\bibitem{Toth1971}
J.~T\'oth,
``State equations of the solid--gas interface layers,''
\textit{Acta Chim.\ Acad.\ Sci.\ Hung.}\ \textbf{69}, 311--328 (1971).

\bibitem{Jaroniec1975}
M.~Jaroniec,
``Adsorption on heterogeneous surfaces: the exponential equation
for the overall adsorption isotherm,''
\textit{Surface Sci.}\ \textbf{50}, 553--564 (1975).

\bibitem{Redlich1959}
O.~Redlich and D.L.~Peterson,
``A useful adsorption isotherm,''
\textit{J.\ Phys.\ Chem.}\ \textbf{63}, 1024 (1959).

\bibitem{BET1938}
S.~Brunauer, P.H.~Emmett, and E.~Teller,
``Adsorption of gases in multimolecular layers,''
\textit{J.\ Am.\ Chem.\ Soc.}\ \textbf{60}, 309--319 (1938).

\bibitem{Vieth1976}
W.R.~Vieth, J.M.~Howell, and J.H.~Hsieh,
``Dual sorption theory,''
\textit{J.\ Membr.\ Sci.}\ \textbf{1}, 177--220 (1976).

\bibitem{Nederlof1990}
M.M.~Nederlof, W.H.~van Riemsdijk, and L.K.~Koopal,
``Determination of adsorption affinity distributions,''
\textit{J.\ Colloid Interface Sci.}\ \textbf{135}, 410--426 (1990).

\bibitem{Cernik1995}
M.~Cernik, M.~Borkovec, and J.C.~Westall,
``Regularized least-squares methods for the calculation of discrete
and continuous affinity distributions for heterogeneous sorbents,''
\textit{Environ.\ Sci.\ Technol.}\ \textbf{29}, 413--425 (1995).

\bibitem{Jovanovic1969}
D.S.~Jovanovic,
``Physical adsorption of gases. I: Isotherms for monolayer and
multilayer adsorption,''
\textit{Kolloid-Z.\ u.\ Z.\ Polymere}\ \textbf{235}, 1203--1213 (1969).

\bibitem{Dubinin1960}
M.M.~Dubinin and L.V.~Radushkevich,
``The equation of the characteristic curve of activated charcoal,''
\textit{Proc.\ Acad.\ Sci.\ USSR Phys.\ Chem.\ Sect.}\
\textbf{55}, 331--333 (1947).

\bibitem{DubininAstakhov1971}
M.M.~Dubinin and V.A.~Astakhov,
``Development of the concepts of volume filling of micropores in the
adsorption of gases and vapors by microporous adsorbents,''
\textit{Bull.\ Acad.\ Sci.\ USSR Div.\ Chem.\ Sci.}\
\textbf{20}, 3--7 (1971).

\bibitem{Hill1910}
A.V.~Hill,
``The possible effects of the aggregation of the molecules of
haemoglobin on its dissociation curves,''
\textit{J.\ Physiol.}\ \textbf{40}, iv--vii (1910).

\end{thebibliography}

\end{document}
